% ---------------------------------------------------------------------------
% header.tex — opmaak voor de PDF-versie van het paper
%
% Engine: pdflatex. Dat is bewust: op deze machine staat alleen
% texlive-base/latex-*, geen texlive-luatex of -xetex. pdflatex kan geen
% Unicode lezen, dus elk niet-ASCII teken uit de tekst krijgt hieronder een
% expliciete definitie (zie \newunicodechar). De lijst is gemaakt door het
% paper te scannen, niet uit het hoofd — staat er ooit een nieuw symbool in,
% dan meldt pdflatex dat als "Unicode character not set up".
%
% Lettertypen: Palatino (mathpazo, via pandoc -V fontfamily) voor tekst en
% wiskunde, TeX Gyre Heros voor koppen, Latin Modern Mono voor code.
% ---------------------------------------------------------------------------

\usepackage[scaled=0.92]{helvet}   % tgheros kent de optie 'scaled' niet
\renewcommand{\ttdefault}{lmtt}
\usepackage{amsmath}
\usepackage{newunicodechar}

% ── niet-ASCII tekens ──────────────────────────────────────────────────────
\newunicodechar{δ}{\ensuremath{\delta}}
\newunicodechar{τ}{\ensuremath{\tau}}
\newunicodechar{θ}{\ensuremath{\theta}}
\newunicodechar{κ}{\ensuremath{\kappa}}
\newunicodechar{λ}{\ensuremath{\lambda}}
\newunicodechar{Δ}{\ensuremath{\Delta}}
\newunicodechar{−}{\textminus}
\newunicodechar{≤}{\ensuremath{\leq}}
\newunicodechar{√}{\ensuremath{\surd}}
\newunicodechar{→}{\ensuremath{\rightarrow}}
% δ met punt erboven (stuursnelheid): het bouwscript vervangt de combinatie
% δ + U+0307 door dit ene teken, zodat het ook binnen \texttt klopt.
\newunicodechar{}{\ensuremath{\dot{\delta}}}

\usepackage{xcolor}
\definecolor{inkblue}{HTML}{2F6FD8}
\definecolor{inkorange}{HTML}{C4621F}
\definecolor{inkgrey}{HTML}{5B6676}
\definecolor{inkline}{HTML}{C8D0DA}
\definecolor{boxbg}{HTML}{F4F6F9}

\usepackage{titlesec}
\titleformat{\section}{\sffamily\large\bfseries}{}{0pt}{}
\titleformat{\subsection}{\sffamily\normalsize\bfseries}{}{0pt}{}
\titlespacing*{\section}{0pt}{1.6em plus .3em minus .2em}{0.5em}
\titlespacing*{\subsection}{0pt}{1.1em plus .2em minus .1em}{0.35em}

\usepackage{fancyhdr}
\newcommand{\repourl}{https://github.com/chat-l18l/no_cutting_corners_with_AI}
\newcommand{\repotext}{github.com/chat-l18l/no\_cutting\_corners\_with\_AI}
\newcommand{\paperfoot}{%
  \fancyhf{}%
  \fancyfoot[L]{\footnotesize\sffamily\color{inkgrey}No Cutting Corners}%
  \fancyfoot[C]{\footnotesize\sffamily\href{\repourl}{\repotext}}%
  \fancyfoot[R]{\footnotesize\sffamily\color{inkgrey}\thepage}%
  \renewcommand{\headrulewidth}{0pt}%
  \renewcommand{\footrulewidth}{0.4pt}}
\pagestyle{fancy}\paperfoot
\fancypagestyle{plain}{\paperfoot}   % ook op de titelpagina

% ── kader met repository-verwijzing en QR-code, onder de titel ─────────────
\usepackage{graphicx}
\newcommand{\repobox}{%
  \begin{center}%
  \setlength{\fboxsep}{11pt}%
  \fcolorbox{inkline}{boxbg}{%
    \begin{minipage}{0.88\linewidth}%
      \begin{minipage}[c]{0.70\linewidth}%
        {\sffamily\bfseries\color{inkorange}Repository}\\[4pt]
        {\small Everything in this report --- the C planner, its host tests,
         the plan documents, the GUI and the scripts that produced these
         figures --- is public, for colleagues and students to read, fork and
         use.}\\[8pt]
        {\sffamily\small\href{\repourl}{\repotext}}
      \end{minipage}\hfill
      \begin{minipage}[c]{0.24\linewidth}%
        \includegraphics[width=\linewidth]{img/qr_repo.png}
      \end{minipage}%
    \end{minipage}}%
  \end{center}\par\medskip}

% ── brede figuur op een liggende pagina ────────────────────────────────────
% lscape verwisselt \textwidth en \textheight binnen de omgeving, dus
% \linewidth is hier al de volle breedte van de gedraaide pagina (~25 cm).
% Het onderschrift staat smaller: een regel van 25 cm leest niet.
\usepackage{pdflscape}
\newcommand{\widefigopen}[1]{%
  \begin{landscape}%
  \vspace*{\fill}%
  \begin{center}\includegraphics[width=\linewidth]{#1}\end{center}%
  \par\medskip
  \begin{center}\begin{minipage}{0.86\linewidth}\small}
\newcommand{\widefigclose}{%
  \end{minipage}\end{center}\par\vspace*{\fill}\end{landscape}}

% Onderschriften zijn gewone alinea's (geen floats): figuur en tekst blijven
% dus staan waar ze in de bron staan, wat voor een rapport prettiger leest.
\setlength{\emergencystretch}{3em}
